\emph{Översikt}
Listan, tupeln och uppslagslistan täcker det mesta --- men inte allt.
Den här modulen inför tre behållare till, och de tre svarar på var sin
fråga.  \emph{Mängden} svarar på \enquote{vilka olika finns här?}: den
har inga dubbletter och ingen ordning, och kan jämföras med andra
mängder med union, snitt och differens.  \emph{Stacken} och \emph{kön}
svarar på \enquote{vilket element ska jag ta härnäst?}, och de svarar
olika: stacken lämnar ut det som lades in sist, kön det som lades in
först.  Vi bygger stacken av en lista och kön av
\mintinline{python}{collections.deque}, och ser efter i Pythons
dokumentation varför det inte blir tvärtom.  Modulen avslutas med att
sätta alla behållarna bredvid varandra: valet av behållare styrs av
vilka operationer programmet behöver.

\emph{Lärandemål}
Efter denna modul ska studenten kunna:
\begin{restatable}{lo}{ContLOSet}\label{ContLOSet}%
  Skapa och använda en mängd --- unika element, medlemskap, samt union,
  snitt och differens --- och förklara varför en mängd varken har
  ordning eller dubbletter.
\end{restatable}% Lo09
\begin{restatable}{lo}{ContLOStack}\label{ContLOStack}%
  Använda en lista som stack, med \mintinline{python}{append} och
  \mintinline{python}{pop}, och förklara sist-in-först-ut.
\end{restatable}% Lo09
\begin{restatable}{lo}{ContLOQueue}\label{ContLOQueue}%
  Använda \mintinline{python}{collections.deque} som kö, med
  \mintinline{python}{append} och \mintinline{python}{popleft}, förklara
  först-in-först-ut, och motivera varför en lista är ett sämre val för
  en kö.
\end{restatable}% Lo09
\begin{restatable}{lo}{ContLOChoose}\label{ContLOChoose}%
  Välja behållare --- lista, tupel, mängd, uppslagslista, stack eller
  kö --- utifrån vilka operationer programmet behöver, och motivera
  valet.
\end{restatable}% Lo09, Lo02
\begin{restatable}{lo}{ContLODocs}\label{ContLODocs}%
  Hitta och läsa dokumentationen för Pythons behållare och därifrån
  använda en behållare eller metod som inte gåtts igenom på
  föreläsningen.
\end{restatable}% Lo09

\emph{Förkunskaper}
Modulen förutsätter listor, tupler och uppslagslistor, samt
\mintinline{python}{for}- och \mintinline{python}{while}-slingor och
egna funktioner.

\emph{Läsning}
Kursens läsförståelseövning \emph{Dokumentation om olika behållare}
arbetar med samma dokumentation som modulen hänvisar till.
